\documentclass[10.5pt]{article}
\usepackage[margin=0.75in]{geometry}
\usepackage{amsmath}
\usepackage{booktabs}
\usepackage{array}
\usepackage{enumitem}
\usepackage{titlesec}
\usepackage{parskip}
\usepackage[hidelinks]{hyperref}
\usepackage{microtype}

\titlespacing*{\section}{0pt}{6pt}{3pt}
\titlespacing*{\subsection}{0pt}{4pt}{1pt}
\titleformat{\section}{\normalfont\large\bfseries}{\thesection}{0.6em}{}
\titleformat{\subsection}{\normalfont\bfseries}{\thesubsection}{0.6em}{}

\setlist{itemsep=0pt, topsep=1pt, parsep=0pt}
\setlength{\parskip}{1pt}
\setlength{\tabcolsep}{4pt}
\renewcommand{\arraystretch}{0.88}
\pagestyle{empty}

\title{\vspace{-1.5em}Design Note --- Lexical vs.\ Semantic Retrieval for News Recommendation\vspace{-0.5em}}
\author{\small CS4.406 Information Retrieval \& Extraction --- Assignment 1}
\date{}

\begin{document}
\maketitle
\vspace{-2em}

\section{What We Built}

News recommendation differs structurally from catalog domains like retail
or streaming: item relevance half-life is hours, not months, so item
cold-start is the \emph{default} state, not an edge case (Wu et al.\ 2020;
Kruse et al.\ 2024) --- motivating content-based candidate retrieval,
lexical (BM25) and semantic (embeddings), over collaborative filtering.

This project implements both methods, plus one shared evaluation harness,
over MIND and EB-NeRD under a unified schema and temporal-split protocol
(ADR-001, ADR-002), comparing methods and datasets on identical footing
rather than via method-specific pipelines. Both were pushed through to
real Codabench leaderboard submissions on both competitions (\S3.5). Each
dataset's weakest link was then chased with a follow-on candidate ---
NRMS-lite for MIND (\S3.6), a LightGBM ranker for EB-NeRD (\S3.7) --- both
of which became this project's real leaderboard wins.

\section{Choices We Made}

\textbf{Data \& splits.} Official train/validation splits are used as-is
(7-day windows) --- the only feasible option given data budgets (ADR-001).
A unified three-table schema (articles, impressions, user\_history) with
mandatory core + dataset-specific optional fields lets one retrieval and
evaluation code path run against both datasets (ADR-002).

\textbf{Lexical retrieval.} BM25Okapi over unweighted concatenated
title+abstract query text (full click history, stopwords removed), scored
as a precomputed sparse weight matrix rather than \texttt{rank\_bm25}'s
naive \texttt{get\_scores()} --- the latter measured at $\sim$10hr projected
vs.\ $\sim$80s actual at MINDsmall-dev scale, a real infeasibility, not a
style preference (ADR-005, ADR-006).

\begin{sloppypar}
\textbf{Semantic retrieval.} One embedding model computed over both datasets
(rejecting EB-NeRD's provided embeddings + a separate MIND model, to keep
the single-code-path comparison valid) --- \texttt{paraphrase-multilingual-MiniLM-L12-v2},
chosen over \texttt{multilingual-e5-small} after a
same-category vs.\ different-category cosine-similarity discrimination
check: MiniLM's gap (0.069, 2.4x ratio) is real and usable; e5-small's
\texttt{passage:}-prefixed convention compresses nearly all pairs into a
narrow 0.75--0.78 band (gap 0.017, ratio 1.02x) --- e5's asymmetric
query/passage objective doesn't fit this project's symmetric
user-profile-vs-article use case (ADR-008). ANN backend is brute-force
cosine similarity --- 0.99ms/query at MINDsmall scale, 2.44ms/query at
MINDlarge's 255{,}990 users --- FAISS stays unjustified at this scale
(ADR-008 + Addendum; revisited at 10x scale, \S5). User representation is
mean-pooled, L2-renormalized embeddings of full click history; zero-history
users get an explicit all-zero vector, mirroring BM25's own cold-start
reporting posture.
\end{sloppypar}

Rejecting EB-NeRD's provided embeddings (above) was argued, not measured,
until a follow-up (ADR-008 Addendum) found the provided artifact wins
recall@K/accuracy CI-clear (AUC 0.5453 vs.\ 0.5430) but loses Diversity@10
--- a real, EB-NeRD-only trade-off, not a blanket verdict.

\textbf{Evaluation.} Two distinct questions are measured, not conflated:
\emph{recall@K} (does retrieval surface the clicked article out of the
\emph{whole} corpus, $K \in \{50,100,200\}$) and \emph{Q4's ranking
metrics} --- AUC, MRR, nDCG@5, nDCG@10, intra-list diversity, novelty,
catalog coverage --- over candidates \emph{already listed in each
impression}. Ties break via a deterministic per-impression-seeded
pseudo-random permutation (cold users produce exact ties); $K=10$ for
diversity/novelty. Every metric reports 95\% bootstrap CIs ($n=2000$),
sliced warm/cold (history length $\geq 5$ = warm), except coverage, a
point estimate only --- set-union statistics are mechanically biased
under with-replacement resampling (ADR-007).

\section{Observations}

\subsection*{3.1 Recall@K --- candidate generation over the whole corpus}

{\footnotesize
\begin{tabular}{@{}llccc@{}}
\toprule
Dataset (split) & Method & recall@50 & recall@100 & recall@200 \\
\midrule
MIND-small dev, overall & BM25 & 0.73\% & 1.50\% & 2.62\% \\
MIND-small dev, overall & Embed & 0.90\% & 1.60\% & 2.78\% \\
MIND-small dev, warm & BM25 & 0.73\% & 1.54\% & 2.73\% \\
MIND-small dev, warm & Embed & 0.93\% & 1.67\% & 2.92\% \\
MIND-small dev, cold & BM25 & 0.72\% & 1.19\% & 1.78\% \\
MIND-small dev, cold & Embed & 0.65\% & 1.10\% & 1.81\% \\
EB-NeRD-demo val (warm only*) & BM25 & 1.01\% & 2.13\% & 4.02\% \\
EB-NeRD-demo val (warm only*) & Embed & 0.32\% & 0.94\% & 2.57\% \\
EB-NeRD-small val (warm only*) & BM25 & 0.72\% & 1.44\% & 2.77\% \\
EB-NeRD-small val & Embed & 0.14\% & 0.43\% & 1.21\% \\
\bottomrule
\end{tabular}\par}

{\footnotesize *EB-NeRD's active-user-filtered bundles have zero cold-start
users (min history length = 5) by construction. Source:
\texttt{experiments/\allowbreak\{bm25,embed\}\_\allowbreak\{mind,ebnerd,ebnerd\_small\}\_\allowbreak2026-08-10/}.
Not run at MINDlarge scale --- flagged in \S6, not approximated.\par}

\subsection*{3.2 Q4 ranking metrics --- AUC / MRR / nDCG (bootstrap 95\% CI)}

{\footnotesize
\begin{tabular}{@{}llccc@{}}
\toprule
Dataset & Method & AUC overall (CI) & AUC warm & AUC cold \\
\midrule
MIND-small dev & BM25 & 0.569 (0.567--0.571) & 0.577 & 0.524 \\
MIND-small dev & Embed & 0.634 (0.632--0.636) & 0.644 & 0.574 \\
MINDlarge dev & BM25 & 0.570 (0.569--0.571) & 0.578 & 0.522 \\
MINDlarge dev & Embed & 0.634 (0.633--0.634) & 0.643 & 0.575 \\
EB-NeRD-demo val & BM25 & 0.533 (0.528--0.537) & --- (=overall) & n/a \\
EB-NeRD-demo val & Embed & 0.544 (0.539--0.549) & --- & n/a \\
EB-NeRD-small val & BM25 & 0.529 (0.527--0.530) & --- & n/a \\
EB-NeRD-small val & Embed & 0.543 (0.541--0.544) & --- & n/a \\
\bottomrule
\end{tabular}\par}

On MIND (both scales), embeddings' AUC CI does not overlap BM25's on
either cohort, and widens the relative warm/cold disparity rather than
closing it (embeddings lift warm AUC $\sim$0.065--0.067 over BM25 but
cold by only $\sim$0.050--0.053). On EB-NeRD, embeddings' AUC CI also
doesn't overlap BM25's --- a modest ($\sim$0.011--0.014) edge --- despite
BM25 winning recall@K by 1.6--2.3x on the \emph{same} corpus (\S3.1):
recall@K asks whether the click is anywhere in the whole catalog, AUC
asks how well the method ranks EB-NeRD's already-curated, shorter
candidate lists (median 9--12 vs.\ MIND's 23) --- see \S6. MRR/nDCG@5/@10
follow AUC's direction on both datasets (e.g.\ MINDlarge nDCG@10: 0.390
vs.\ 0.350; EB-NeRD-small: 0.459 vs.\ 0.454).

\subsection*{3.3 Diversity / novelty / coverage (K=10, point estimate for coverage)}

{\footnotesize
\begin{tabular}{@{}llccc@{}}
\toprule
Dataset & Method & Diversity@10 & Novelty@10 & Coverage@10 \\
\midrule
MIND-small dev & BM25 & 0.837 & 16.30 & 0.083 \\
MIND-small dev & Embed & 0.827 & 16.15 & 0.078 \\
MINDlarge dev & BM25 & 0.836 & 18.09 & 0.067 \\
MINDlarge dev & Embed & 0.828 & 17.76 & 0.065 \\
EB-NeRD-small val & BM25 & 0.795 & 17.17 & 0.206 \\
EB-NeRD-small val & Embed & 0.789 & 17.19 & 0.205 \\
\bottomrule
\end{tabular}\par}

BM25 and embeddings sit within $\sim$1 point of each other on every
beyond-accuracy metric, on every corpus --- accuracy differs meaningfully,
beyond-accuracy behavior barely does ($n$ up to 376{,}471 impressions, so
not automatically noise, but not the large effect AUC shows).

\subsection*{3.4 System cost}

BM25 at MINDlarge scale: 1.70s index build, 23.5MB sparse weight matrix,
$\sim$9.8 min projected full retrieval for 255{,}990 users (ADR-006
Addendum). Embeddings: 264.4s to encode+cache all 72{,}023 articles,
$\sim$10.4 min projected retrieval (ADR-008 Addendum). Both stay local on
an 8GB machine at this scale --- only EB-NeRD's blind test split (13.5M
impressions, no local ground truth) required Kaggle.

\subsection*{3.5 Leaderboard results (Q5)}
\vspace{-2pt}

{\footnotesize
\begin{tabular}{@{}llccc@{}}
\toprule
Competition & Method & Score & Submission ID & Date \\
\midrule
MIND (competitions/13967) & Embed (MiniLM) & 0.6195 & 886468 & 2026-08-12 13:01 \\
MIND (competitions/13967) & Cohort-gated combiner & 0.6192 & 896696 & 2026-08-22 \\
MIND (competitions/13967) & \textbf{NRMS-lite (Cand.\ J)} & \textbf{0.6462} & 901961 & 2026-08-26 \\
EB-NeRD (competitions/2469) & Embed (MiniLM) & 0.5404 & 888045 & 2026-08-13 23:08 \\
EB-NeRD (competitions/2469) & Contrastive vector & 0.5404 & 896072 & 2026-08-21 13:01 \\
EB-NeRD (competitions/2469) & \textbf{GBDT ranker (Cand.\ K, \S3.7)} & \textbf{0.7542} & 907863 & 2026-08-29 23:01 \\
\bottomrule
\end{tabular}\par}
\vspace{-3pt}

Both used the embedding method --- the clear winner on local dev-set AUC
for both datasets (\S3.2). EB-NeRD's first score (0.5404) is close to
local validation AUC on \texttt{ebnerd\_small} (0.5430) --- a coherence
check, not a claim the two are defined identically.

\textbf{Second MIND submission --- cohort-gated combiner, essentially
flat.} Local validation predicted a CI-clear win at every stage
(+0.0027 MINDsmall-dev, +0.0019 MINDlarge-dev re-verification, ADR-010's
Addendum), but the real result was \textbf{0.6192 vs.\ 886468's 0.6195
(-0.0003)} --- flat, not the predicted win. No pipeline defect found
(line count, malformed-permutation, all 32 \texttt{N89741} cases clean);
most likely a genuine population difference, since MIND's dev/test splits
are disjoint calendar weeks (ADR-001) and this edge was built around
MINDlarge-dev-specific properties.

\textbf{Second EB-NeRD submission --- contrastive vector, real test set.}
Both submissions round to the same Score (0.5404) --- checked, not
assumed identical: different checksums, 99.48\% of impressions rank
differently for the same ID, confirming independent runs. Codabench's
per-submission detail view (grouped by date, ``50\% of the testset'')
resolves the tie: contrastive leads 5/8 days and the mean AUC (0.5402
vs.\ 0.5397, +0.0005, sign flips three times day-to-day) --- real but far
smaller than the \textbf{$\sim$0.0023 CI-clear gap} ADR-008's Addendum
measured locally (0.5453 vs.\ 0.5430).

\textbf{Third MIND submission --- NRMS-lite (Candidate J), a real win.}
A trainable title/history self-attention architecture (Wu et al.\ 2019,
GloVe-init, ADR-012), unlike every frozen-embedding candidate below.
Trail: MINDsmall-dev 0.6391 $\to$ MINDlarge-dev 0.6579 (margin
\emph{widened}, unique among all candidates) $\to$ real leaderboard
\textbf{0.6462} (\textbf{+0.0267} over the original submission ---
compressed from the dev screen, but a real win, unlike the combiner above).
A narrow-impact scoring-catalog bug (initially estimated 32/2,370,727
impressions, later measured directly at 2,087/2,370,727, 0.088\%) was found
and fixed post-submission; a corrected resubmission scored an identical
\textbf{0.6462}, confirming the bug never put the result in question
(ADR-012).

\subsection*{3.6 MIND candidate search: ten approaches, and what worked}

Ten approaches screened on MINDsmall-dev vs.\ baseline (0.634, \S3.2)
before J, all atop \emph{frozen} embeddings with a shallow head. Only
G's routing won CI-clear overall (didn't hold at real scale), only F's
combiner won CI-clear on the cold cohort. Evidence: ADR-010 (D--H)/011 (I)/012 (J).

{\tiny
\renewcommand{\arraystretch}{0.75}
\begin{tabular}{@{}llc@{}}
\toprule
Cand.\ & Method & Result \\
\midrule
A & Entity embeddings (TransE) & 0.553, loss \\
B & 50/50 BM25+embedding hybrid & 0.626, loss \\
C & Recency-weighted embedding query & 0.627, loss \\
D & Symbolic category/entity overlap & 0.613, loss \\
E & Train-popularity only & 0.532, loss \\
F & 7-feature LR combiner & 0.626 loss; 0.593 cold win \\
G & Cohort-gated routing & 0.637 local win; 0.6192 real (flat) \\
H1 & LogisticRegression, balanced & 0.632, loss \\
H2 & HistGradientBoosting & 0.599, loss \\
I & Attention re-ranker, 5 epochs & 0.623, loss \\
I-long & Same, 30 epochs & 0.621, worse than I \\
I-pop & I-long + raw popularity feature & 0.513, inconclusive \\
\textbf{J} & \textbf{Trainable title+history encoders} & \textbf{0.658 local win; 0.6462 real win} \\
\bottomrule
\end{tabular}\par}

\subsection*{3.7 EB-NeRD candidate search: Candidate K (GBDT learning-to-rank)}

EB-NeRD's deployed score (0.5404) sits \emph{below} the challenge's own
popularity baseline (0.5970): in-view candidates are drawn from the same
front page at the same moment (median age 3.7h, 58\% under 6h), leaving
content similarity little to discriminate on. Unlike MIND's ten-way search
(\S3.6), EB-NeRD's search is a single candidate, K --- a LightGBM
learning-to-rank model over 65 engineered behavioural/temporal features
--- chosen because the RecSys Challenge 2024 organizers' own report states
most real competitors used GBDT ensembles, not neural architectures; a
better-embeddings option and an NRMS-style option (as Candidate J did for
MIND) were both rejected on that evidence before building K (ADR-013).

\begin{sloppypar}
\textbf{A real per-impression bug was found and fixed mid-session.}
Short-term (recency-decayed) features were computed once per \emph{user}
against a shared reference time, not once per \emph{impression} ---
effectively static. Fixed (\texttt{compute\_short\_term\_features},
referenced to each impression's own timestamp); a second, independently
found leak-safety gap in the same path (no upper bound excluding a
future-dated history click) was closed at the same time. Local
\texttt{ebnerd\_small} AUC moved from \textbf{0.7514--0.7528 to
0.7581--0.7597}, CI-clear (ADR-013's 2026-08-27 Addendum).
\end{sloppypar}

\textbf{Trained at real \texttt{ebnerd\_large} scale on Ada} (IIIT-H's
SLURM cluster): 12{,}063{,}890 train / 12{,}566{,}385 validation
impressions, the latter scored in full via a streaming path. Result: local
validation AUC \textbf{0.7590 (95\% CI 0.7588--0.7593)} --- essentially
unmoved by $\sim$10.7x more training data, a real null result on scale for
this feature set (ADR-013's 2026-08-29 Addendum).

\textbf{Real Codabench result: 0.7542} (\S3.5) --- \textbf{+0.2138} over
the previous deployed EB-NeRD submission, \textbf{+0.1572} over the
challenge's popularity baseline. Compression from local
\texttt{ebnerd\_large} validation was small (0.7590 $\to$ 0.7542,
$-$0.0048) --- the first EB-NeRD/MIND candidate here whose CI-clear local
win transferred largely intact, rather than evaporating (MIND's
cohort-gated combiner) or thinning to near-nothing (EB-NeRD's contrastive
vector), plausibly because its edge is structural, leak-safe item-level
signal rather than a validation-population-specific pattern (ADR-013's
2026-08-30 Addendum).

Feature importance at real scale: freshness (led by \texttt{article\_age\_h})
is \textbf{33.3\%} of gain; long-term interest 14.1\%; \textbf{short-term
interest only 0.8\%} (down from an already-small 1.9\% pre-fix at
small scale). The BlackPearl-derived long/short-term hypothesis this
candidate was built to test is therefore \textbf{falsified}: what wins is
freshness \emph{relative to other candidates in the same list} and
\texttt{context\_embed\_sim} (what the user is reading right now), not
modelled user interest. Withholding position features \emph{improved} AUC
by a CI-clear +0.0011 --- the shipped arm (\texttt{K\_rank\_nopos}) omits
them and needs no such caveat.

\section{Anti-Gaming and Leakage (Q9)}

\begin{sloppypar}
\textbf{Serving-time-unavailable features --- measured, not just avoided.}
Novelty's popularity signal comes from the \textbf{train split only},
Laplace-smoothed, never the evaluation split (ADR-007). EB-NeRD's raw
\texttt{articles.parquet} also carries three lifetime aggregate fields
--- \texttt{total\_inviews}, \texttt{total\_pageviews},
\texttt{total\_read\_time} --- unavailable at serving time;
\texttt{src/datasets/ebnerd.py} never reads them (grep-verified).
\texttt{scripts/run\_leakage\_ablation.py} measures what including them
would buy: an untuned 50/50 blend of the deployed embedding score with a
min-max-normalized leak signal, on \texttt{ebnerd\_small} validation
(ADR-009):
\end{sloppypar}

{\footnotesize
\begin{tabular}{@{}lccc@{}}
\toprule
Metric & Without leaky features & With leaky features (50/50 blend) & CIs overlap? \\
\midrule
AUC & 0.5430 (0.5415--0.5445) & 0.5662 (0.5646--0.5678) & No \\
MRR & 0.3437 (0.3421--0.3452) & 0.3638 (0.3623--0.3653) & No \\
nDCG@5 & 0.3804 (0.3784--0.3823) & 0.4034 (0.4016--0.4051) & No \\
nDCG@10 & 0.4591 (0.4574--0.4607) & 0.4783 (0.4767--0.4799) & No \\
\bottomrule
\end{tabular}\par}

Every metric moves by a CI-clear margin from exposure to these fields
alone (AUC +0.023 absolute, $\sim$4.3\%), untuned. $\sim$48\% of articles
were median-imputed, not zero-filled, so this gap is plausibly a
conservative underestimate --- direct, measured evidence for excluding
these fields entirely (ADR-009).

\textbf{Leakage-boundary test.} \texttt{tests/integration/test\_leakage.py}
asserts no user's history click occurs at or after their own earliest
impression time. \textbf{Exists and passes for EB-NeRD} (demo train +
validation splits). \textbf{Explicitly skip-marked for MIND}: MIND
provides only a static per-user article-ID list, no per-click timestamps,
so the temporal boundary can't be independently verified from the data.
An acknowledged gap, not an oversight --- MIND's temporal safety instead
rests on the dataset's own documented split construction (ADR-001).

\section{Where It Breaks at 10$\times$}

Every real scaling failure so far was invisible at the previously-exercised
scale and appeared at a 5--14x jump, not a 2x one --- a pattern, not a
one-off.

\textbf{MINDsmall $\rightarrow$ MINDlarge} ($\sim$5x users: 50{,}000
$\rightarrow$ 255{,}990; $\sim$14x impression-candidate rows) hung in a
per-token Python loop never exercised at MINDsmall's counts; the first
fix projected to \textbf{$\sim$27GB} (measured, not estimated) from one
Python object per repeated ID string, resolved via categorical dtype
(12.5x reduction, $\sim$2.2GB); a second bottleneck
(\texttt{series.map()} on categorical data, $>$1hr, zero progress) was
found only by profiling the stuck process. \textbf{MINDlarge\_test (2.37M)
$\rightarrow$ ebnerd\_testset (13.5M, 5.7x further)} hit a third,
structurally identical failure: a converter materializing a full
\texttt{list[dict]} before writing projected to $\sim$16GB, invisible at
the smaller scale --- fixed by a generator consuming one row at a time.
\textbf{\texttt{ebnerd\_small} $\rightarrow$ real \texttt{ebnerd\_large}
training on Ada} ($\sim$52x further train impressions) hit a fourth,
real SLURM-confirmed OOM: an un-split 65-feature train matrix
($\sim$11.5GB) was never freed, and a column-sliced copy needed by one arm
coexisted with it ($\sim$22.7GB from four arrays alone). Fixed by freeing
the matrix once training finished, capping the training sample
4M$\to$2.5M impressions, and raising \texttt{--mem-per-cpu} 4G$\to$8G
(four further non-modelling Ada incidents the same session in ADR-013's
2026-08-29 Addendum).

Projecting one more 10x from the largest verified scale is inference, not
measurement: (1) retrieval time scales roughly linearly with users, so
10x users alone stays locally feasible ($\sim$100 min); (2) brute-force
cosine is $O(\text{users}\times\text{articles})$, so a 10x corpus
alongside 10x users is 100x compute ($\sim$10.4 min $\to \sim$17 hours),
forcing a revisit of ADR-008's ``FAISS unjustified'' call; (3) every real
bug so far shared one shape --- a per-row-object or full-materialization
pattern past an unexercised scale --- and nothing guarantees another path
is safe at a further 10x.

\section{Limitations and Open Questions}
\enlargethispage{3\baselineskip}

\begin{sloppypar}
\begin{itemize}
\item \textbf{EB-NeRD's real result (0.7542, \S3.7) did not reach this
  project's internally-discussed 0.80 target --- stated plainly.} The
  target referenced a leakage-dependent figure: the RecSys Challenge 2024
  winner scored 88.64 in the organizers' own ablation arm, but
  \textbf{76.99} once those same features were found to leak future
  information and removed (an 11.65-point drop; arXiv:2409.20483). 76.99
  is the honest ceiling this project can defensibly compare against;
  0.7542 sits \textbf{0.0157 below it}, on
  features that passed the same leakage audit every candidate here uses.
\item \textbf{True (zero-history) cold-start is a shared ceiling}, not
  solved by either method (MIND cold-cohort recall@200: BM25 1.78\% vs.\
  embed 1.81\%); EB-NeRD's cold cohort is structurally empty in both
  \texttt{demo}/\texttt{small} (min history length = 5, \texttt{large}
  unverified), and recall@K (\S3.1) was never run at MINDlarge scale.
\item \textbf{The Q9 leakage-boundary test is real but partial} (passes
  EB-NeRD, skip-marked MIND, no per-click timestamps; \S4); its ablation
  magnitude is a lower bound, not a maximum, from an untuned 50/50 blend.
  Two other choices rest on a proxy, not a gold standard: the encoder's
  discrimination check (ADR-008) and the stopword list (ADR-005).
\item \textbf{The EB-NeRD recall-vs-AUC disagreement isn't fully
  explained} (\S3.2): BM25 wins whole-corpus recall 1.6--2.3x, yet
  embeddings edge the curated candidate-list AUC --- plausibly a diffuse
  mean-pooled query, not isolated by experiment.
\item \textbf{Local wins compress at real scale, but don't always
  vanish}: MIND's combiner and EB-NeRD's contrastive vector compressed to
  flat; Candidate J widened at MINDlarge-dev, then compressed at the real
  leaderboard but stayed substantial (+0.0267 over baseline, \S3.5/3.6);
  \textbf{Candidate K compressed least of all} (0.7590 local $\to$ 0.7542
  real, $-$0.0048, \S3.7) --- architecture/signal type plausibly affects
  compression (inference, $N=4$). Relatedly, the BlackPearl long/short-term
  hypothesis (A1, ADR-013) is \textbf{falsified}: short-term interest fell
  to 0.8\% of Candidate K's gain \emph{after} the recency fix made it more
  correctly computed, not less important.
\item \textbf{A numpy/CPython 3.14 mismatch produced provably wrong array
  results (ADR-013)}, fixed by rebuilding on Python 3.11 --- unclear if
  any earlier number predates the fix.
\end{itemize}
\end{sloppypar}

\end{document}
